\documentclass[10pt]{article}
\usepackage{resume_shared}

\begin{document}
\fontsize{9.8pt}{11.7pt}\selectfont

\ResumeName{Sri Ravi Kumar Birada}
\ResumeContact{Hyderabad, India}{(+91) 7200999942}{ravikumar31birada@gmail.com}{www.linkedin.com/in/ravi-20043}

\section{Summary}
AI Product Manager with 7+ years across enterprise SaaS and AI products, spanning ideation, build, MVP, evaluation, optimization, and scale. Strong at value stream mapping, single-persona problem framing, outcome-led prioritization, model evaluation, RAG setup, prompt optimization, pilot success definition, and governance-aware AI adoption. Recent work includes a 3-module AI roadmap, agentic AI RFP automation, 0 to 1 SaaS capability shaping, and partner-led pilot design.

\section{Professional Experience}
\RoleEntry{AI Product Manager - Enterprise AI Products}{HCLTech}{2023 -- Present}
\begin{itemize}
  \item Mapped 20+ workflow activities using value stream mapping, translating customer journeys into 5 outcome-led AI use cases across 3 enterprise SaaS product modules and a 12-month roadmap.
  \item Prioritized single-persona, point-based AI solutions by defining user goals, success criteria, MVP scope, and evaluation checkpoints before moving features from ideation to build.
  \item Built an agentic AI RFP response workflow that ingests customer RFPs and enterprise capability / technology-stack data, applies policy and quality validation, and reduces last-minute proposal rework risk.
  \item Established model-evaluation discipline across SOTA model comparison, prompt optimization, RAG setup, task execution, and A/B testing using BERT score, topic adherence, F1, BLEU, and PII-redaction checks.
  \item Optimized GenAI solution readiness through alpha-user testing, RAG chunking experiments, rerun reduction, LLM-cost governance, customer-journey finalization, and fine-tuning assessment where required.
  \item Scaled repeatable AI patterns by assessing connector feasibility, microservices integration, user onboarding, training needs, responsible AI controls, and enterprise deployment readiness.
\end{itemize}

\RoleEntry{Product Manager - Enterprise SaaS Products}{HCLTech}{May 2022 -- 2023}
\begin{itemize}
  \item Captured multi-account customer inputs and converted them into modular SaaS requirements across customer operations, workflow orchestration, reporting, and commercial governance.
  \item Shaped a 0 to 1 enterprise SaaS capability from 6+ customer-account inputs, defining 8 MVP feature areas across platform-extension and standalone packaging paths.
  \item Packaged platform capabilities into adoption narratives by connecting buyer pain points, implementation sequencing, and time-to-value expectations for enterprise users.
  \item Accelerated operational readiness for an enterprise product deployment supporting 300K+ SKUs and 52+ distributed business entities by coordinating execution and reporting enablement, resulting in launch within 21 days.
\end{itemize}

\MiniHeading{Selected Product Outcomes}
\begin{itemize}
  \item Ideation-to-scale model: converted 20+ mapped workflow activities into 5 AI use cases, 3 product-module priorities, and a 12-month execution view.
  \item Evaluation system: formulated golden datasets and metric-backed checks for model output quality, retrieval relevance, topic adherence, and PII handling.
  \item Optimization loop: linked prompt tuning, RAG chunking, rerun reduction, LLM-cost governance, and alpha-user testing to product-readiness decisions.
  \item Partner pilot design: evaluated 25+ partners and onboarded 3 through pilot-account selection, technical evaluation, and 5-metric success definition.
\end{itemize}

\RoleEntry{Network Engineer}{Ericsson Global India Limited}{Oct 2016 -- Jan 2019}
\begin{itemize}
  \item Reduced deliverable turnaround time by 15\% by automating report validation workflows with Python.
  \item Cut server infrastructure requirements by 40\% through performance improvements in analysis tooling and workflow design.
  \item Received the Ericsson Power Award in Q4 2017 for automation impact in Single Site Verification reporting.
\end{itemize}

\section{Education}
\InlineSectionText{MBA, MDI Gurgaon (2022)}
\InlineSectionText{B.Tech, Electronics and Communication Engineering, Amrita School of Engineering (2016)}

\section{Certifications}
\InlineSectionText{OpenAI AI Technical Practitioner (2026) \textbar\ Google Certified GenAI Leader (2025) \textbar\ GenAI Prompt Engineer L2 (2025) \textbar\ Certified Scrum Product Owner (CSPO) (2024)}

\section{Skills}
\SkillGroup{Product Management}{AI Product Management, Product Strategy, Roadmapping, Product-Market Fit, 0 to 1 Module Development, KPI Ownership}
\SkillGroup{Enterprise AI}{GenAI Strategy, Responsible AI, RAG Setup, Prompt Optimization, Agentic Workflow Design, Model Evaluation, LangChain}
\SkillGroup{Transformation and Analytics}{Value Stream Mapping, Workflow Modernization, Partner Onboarding, Pilot Success Definition, A/B Testing, LLM Cost Governance, Power BI, Python}

\end{document}
